\documentclass[11pt,a4paper]{article}
\usepackage{fontspec}
% Load Latin Modern by file, so no system font registration is needed. It
% carries the full Turkish alphabet, dotless i and breves included.
\setmainfont{lmroman10-regular.otf}[
  BoldFont      = lmroman10-bold.otf,
  ItalicFont    = lmroman10-italic.otf,
  BoldItalicFont= lmroman10-bolditalic.otf]
\setsansfont{lmsans10-regular.otf}[
  BoldFont   = lmsans10-bold.otf,
  ItalicFont = lmsans10-oblique.otf]
\setmonofont{lmmono10-regular.otf}[
  BoldFont   = lmmonolt10-bold.otf,
  ItalicFont = lmmono10-italic.otf,
  Scale      = MatchLowercase]
\usepackage[a4paper,margin=2.4cm,bottom=2.6cm]{geometry}
\usepackage{booktabs,longtable,array,fancyvrb,microtype}
\usepackage[table]{xcolor}
\usepackage{enumitem}
\usepackage[hidelinks]{hyperref}
\usepackage{titlesec}
\setlist{nosep,leftmargin=1.3em}
\titleformat{\section}{\large\bfseries}{\thesection}{0.7em}{}
\titleformat{\subsection}{\normalsize\bfseries}{\thesubsection}{0.6em}{}
\titlespacing*{\section}{0pt}{1.4\baselineskip}{0.4\baselineskip}
\titlespacing*{\subsection}{0pt}{1.0\baselineskip}{0.3\baselineskip}
\renewcommand{\arraystretch}{1.15}
\setlength{\parindent}{0pt}
\setlength{\parskip}{0.55\baselineskip}
% The Markdown numbers its sections from 0 (the summary); match it.
\setcounter{section}{-1}
% No Turkish hyphenation patterns in this TeX installation, so give the
% paragraph breaker room instead of letting it overflow the margin.
\sloppy
\emergencystretch=2.5em
\hyphenpenalty=1200
\title{\bfseries Molekülden Kokuya — Teknik Dokümantasyon}
\author{Kamran Heydarov}
\date{15 Ekim 2026}
\begin{document}

\maketitle

\textbf{İnteraktif insan olfaktör reseptör simülatörü} \\
Proje sitesi: \url{https://kammmran.github.io/ENodorIM/#/} \\
Veri sürümü: M2OR \texttt{\small M2OR\_20230428.csv}, derleme tarihi 2026-10-07 \\
Doküman tarihi: 15 Ekim 2026


\vspace{0.4\baselineskip}
\hrule
\vspace{0.9\baselineskip}


\section{Bir bakışta özet}

Bu proje, bir koku molekülünün burna ulaşmasından bilinçli koku algısına kadar olan yolu adım adım gösteren, tarayıcıda çalışan interaktif bir simülatördür. Projenin bilimsel ağırlık merkezi \textbf{reseptör adımıdır}: hangi insan olfaktör reseptörünün (OR) yanıt verdiği, uydurma değerlerden değil, \textbf{yayımlanmış hücre temelli deney verisinden} gelir (M2OR veritabanı). Reseptör adımından sonraki aşamalar (duysal nöron, bulbus olfactorius, korteks) bu aktivasyon örüntüsünün neye yol açtığını gösteren, niteliksel olarak doğru ama \textbf{şematik} modellerdir.

Projenin yöntemsel omurgası tek bir kurala dayanır: \textbf{ekrandaki her sayı, kaynağı ve kanıt düzeyiyle birlikte etiketlenir.} Üç etiket kullanılır:


\begin{small}
\begin{longtable}{@{}p{0.285\linewidth}@{\hspace{5pt}}@{}p{0.645\linewidth}@{}}
\toprule
\textbf{Etiket} & \textbf{Anlamı} \\
\midrule
\endfirsthead
\toprule
\textbf{Etiket} & \textbf{Anlamı} \\
\midrule
\endhead
\bottomrule
\endfoot
\texttt{\small measured} (ölçülmüş) & Yayımlanmış bir deneyden doğrudan gelen değer (örn. EC50) \\
\texttt{\small predicted} (tahmin) & Bir modelin ürettiği değer; modelin kendi doğrulama metriğiyle birlikte gösterilir \\
\texttt{\small schematic} (şematik) & Niteliksel olarak doğru ama parametreleri açıklayıcı amaçlı olan model çıktısı \\
\end{longtable}
\end{small}

Bu ayrım kozmetik değildir: ``test edilmemiş'' ile ``yanıt vermiyor'' sitede \textbf{farklı durumlar} olarak gösterilir ve dışa aktarımda bile korunur. Bu, nörobilimsel yorum için kritik bir ayrımdır — insan OR repertuvarının büyük kısmı için hiçbir ligand bilinmiyor ve bu bilgisizliği sıfır yanıt olarak kodlamak sistematik bir yanlılık üretir.


\vspace{0.4\baselineskip}
\hrule
\vspace{0.9\baselineskip}


\section{Bilimsel çerçeve: üç aşama}

Site, koku algısını literatürdeki standart üç aşamalı kurguyla ele alır:

\begin{enumerate}
\item \textbf{Olfaktör GPCR'ler — uyaranın başlatıcı olayı.} Odorant, olfaktör duysal nöronların silialarındaki OR'lara bağlanır. İnsanda \textasciitilde{}400 işlevsel OR geni vardır ve her nöron (kemirgen verisine göre) tek bir alleli monoallelik olarak ifade eder.
\item \textbf{Sinyalin nöronal entegrasyonu.} Reseptör uyarımı \(\rightarrow\) G\(\alpha\)olf \(\rightarrow\) adenilat siklaz III \(\rightarrow\) cAMP \(\rightarrow\) CNG kanalı \(\rightarrow\) Ca\(^{2+}\) girişi \(\rightarrow\) ANO2 klorür akımı \(\rightarrow\) depolarizasyon \(\rightarrow\) aksiyon potansiyeli dizisi. Aynı reseptör tipini ifade eden nöronlar bulbusta ortak glomerüllere yakınsar; lateral inhibisyon örüntüyü keskinleştirir.
\item \textbf{Sinyalin kod açılımı ve bilinçli algı.} Bulbus çıktısı primer olfaktör kortekse (piriform korteks, AON, tuberculum olfactorium, amigdala, entorhinal korteks) ve oradan OFC, insula, hipokampus, hipotalamus gibi yapılara taşınır.
\end{enumerate}

Simülasyon bu üç aşamayı sırayla canlandırır; ancak \textbf{yalnızca 1. aşama ölçülmüş veriyle sürülür.}


\vspace{0.4\baselineskip}
\hrule
\vspace{0.9\baselineskip}


\section{Veri kaynakları}


\begin{small}
\begin{longtable}{@{}p{0.300\linewidth}@{\hspace{5pt}}@{}p{0.401\linewidth}@{\hspace{5pt}}@{}p{0.229\linewidth}@{}}
\toprule
\textbf{Kaynak} & \textbf{Ne için kullanılıyor} & \textbf{Lisans / atıf} \\
\midrule
\endfirsthead
\toprule
\textbf{Kaynak} & \textbf{Ne için kullanılıyor} & \textbf{Lisans / atıf} \\
\midrule
\endhead
\bottomrule
\endfoot
\textbf{M2OR} (Lalis ve ark., \emph{Nucleic Acids Research} 2024) & Reseptör–odorant çiftleri, EC50 değerleri, tarama sonuçları & Apache-2.0, \href{https://doi.org/10.1093/nar/gkad886}{doi:\allowbreak 10.\allowbreak 1093/\allowbreak nar/\allowbreak gkad886} \\
\textbf{HGNC} & Gen sembollerinin standardizasyonu & \href{https://doi.org/10.1093/nar/gkac888}{doi:\allowbreak 10.\allowbreak 1093/\allowbreak nar/\allowbreak gkac888} \\
\textbf{UniProt} & Protein dizileri, aksesyon eşleme & \href{https://doi.org/10.1093/nar/gkae1010}{doi:\allowbreak 10.\allowbreak 1093/\allowbreak nar/\allowbreak gkae1010} \\
\textbf{GPCRdb} & Jenerik rezidü numaralandırması (Ballesteros–Weinstein benzeri) & \href{https://doi.org/10.1093/nar/gkae1065}{doi:\allowbreak 10.\allowbreak 1093/\allowbreak nar/\allowbreak gkae1065} \\
\textbf{RCSB PDB} & Deneysel kriyo-EM yapıları (örn. 8F76 = OR51E2–propiyonat) & her girdinin birincil atfı \\
\textbf{AlphaFold DB} & Deneysel yapısı olmayan reseptörler için öngörülen modeller (açıkça ``predicted'' etiketli) & \href{https://doi.org/10.1093/nar/gkad1011}{doi:\allowbreak 10.\allowbreak 1093/\allowbreak nar/\allowbreak gkad1011} \\
\textbf{PubChem} & Molekül kimliği, InChIKey, kimyasal yapı & — \\
\textbf{Leffingwell PMP 2001} (Pyrfume üzerinden) & Uzman koku tanımlayıcıları & \textbf{CC BY-NC} \\
\textbf{Dravnieks Atlas} (1985) & Koku karakter profilleri & — \\
\textbf{Keller \& Vosshall} (2016) & İnsan panel değerlendirmeleri & \href{https://doi.org/10.1186/s12868-016-0287-2}{doi:\allowbreak 10.\allowbreak 1186/\allowbreak s12868-016-0287-2} \\
\textbf{Crossref / DataCite} & Her DOI için kaynakça üstverisi & — \\
\end{longtable}
\end{small}

\textbf{Lisans notu.} Leffingwell tanımlayıcı seti CC BY-NC olduğu için sitenin yayımlanmış kopyası \textbf{ticari olmayan} kullanımla sınırlıdır. Akademik ve araştırma kullanımı sorunsuzdur. The Good Scents Company verisi, net bir lisansı olmadığı için \textbf{kullanılmamıştır.}


\vspace{0.4\baselineskip}
\hrule
\vspace{0.9\baselineskip}


\section{Veri işleme hattı ve sayısal durum}

Python hattı (\texttt{\small scripts/pipeline/}) kaynakları indirir, temizler, eşler ve sitenin yüklediği JSON dosyalarını üretir. Her çalıştırmada bir özet raporu basılır (\texttt{\small data/processed/summary.txt}), böylece veri sürümleri arasındaki fark izlenebilir.


\subsection{Adımlar}


\begin{small}
\begin{longtable}{@{}p{0.199\linewidth}@{\hspace{5pt}}@{}p{0.731\linewidth}@{}}
\toprule
\textbf{Modül} & \textbf{Yaptığı iş} \\
\midrule
\endfirsthead
\toprule
\textbf{Modül} & \textbf{Yaptığı iş} \\
\midrule
\endhead
\bottomrule
\endfoot
\texttt{\small sources.py} & Önbelleklenmiş indirmeler (M2OR, HGNC, UniProt, GPCRdb, AlphaFold, PDB, PubChem, Pyrfume, Crossref) \\
\texttt{\small m2or.py} & İnsan tek-bileşenli deneyleri \(\rightarrow\) reseptör–odorant çifti başına \textbf{tek konsensüs sonucu}, M2OR'un yayımlanmış Pair-table kurallarıyla \\
\texttt{\small receptors.py} & M2OR'daki çoklu gen yazımlarını (\texttt{\small hOR10AC}, \texttt{\small OR 10G9}, eski semboller) HGNC sembolüne eşler: sembol \(\rightarrow\) alias \(\rightarrow\) UniProt \(\rightarrow\) dizi \\
\texttt{\small odorants.py} & Test odorantları InChIKey ile; koku tanımlayıcıları \textbf{tam stereokimya eşleşmesiyle} \\
\texttt{\small genotypes.py} & Reseptör genotip vakaları; alleller \textbf{aynı makale içinde} karşılaştırılır \\
\texttt{\small structures.py} & Ligand temasları, yatırılmış PDB koordinatlarından \textbf{ölçülerek} çıkarılır (makaleden alıntılanmaz) \\
\texttt{\small references.py} & Her DOI için atıf üstverisi \\
\end{longtable}
\end{small}


\subsection{Kilit sayılar (mevcut derleme)}


\begin{small}
\begin{longtable}{@{}p{0.391\linewidth}@{\hspace{5pt}}@{}p{0.539\linewidth}@{}}
\toprule
\textbf{Büyüklük} & \textbf{Değer} \\
\midrule
\endfirsthead
\toprule
\textbf{Büyüklük} & \textbf{Değer} \\
\midrule
\endhead
\bottomrule
\endfoot
İnsan tek-bileşenli deney kaydı & 47.316 \\
Konsensüs reseptör–odorant çifti & 37.849 (referans allel 27.503 / varyant allel 10.346) \\
Referans allel çiftleri & agonist \textbf{1.306} \(\cdot\) agonist-değil \textbf{25.827} \(\cdot\) sonuçsuz \textbf{370} \\
Agonistlerin kanıt dağılımı & doz–yanıt (EC50) 545 \(\cdot\) primer tarama 651 \(\cdot\) sekonder tarama 110 \\
Sayısal EC50 taşıyan çift & 513 (referans allel) \(\cdot\) 791 (varyantlar dahil) \(\cdot\) sansürlü (EC50 > x) 32 \\
Sitedeki reseptör ismi & 432 — bunların \textbf{417}'si HGNC'de protein kodlayan OR geni \\
En az bir molekülle test edilmiş reseptör & \textbf{422} \\
En az bir bilinen agonisti olan reseptör & \textbf{235} \\
Tam simülasyonu olan test odorantı & 16 \\
Genotip vakası & 3 (OR7D4, OR2J3, OR10G4) \\
\end{longtable}
\end{small}


\subsection{Veri kalitesi kararları — toplantıda sorulması muhtemel noktalar}

\begin{itemize}
\item \textbf{Çelişen sonuçlar silinmez.} Aynı çift için birbirini tutmayan deneyler varsa çift
\texttt{\small inconclusive} olarak kalır ve ekranda ``sonuçlar uyuşmuyor'' gerekçesiyle gösterilir.

\item \textbf{Medyan EC50.} Birden fazla makale EC50 bildirdiğinde medyan kullanılır, ancak tüm değerler listelenir.
\item \textbf{Tam-log değerler işaretlenir.} Mainland ve ark. 2015'ten gelen 42 referans-allel EC50 değeri tam log basamaklarıdır (±0,5 log) ve gösterildiği her yerde işaretlenir.
\item \textbf{Test aralığı altındaki eğri uydurmaları işaretlenir.} OR51E2 için 9 EC50 değeri 1 nM'nin altına uydurulmuştur, yani bu tahlillerin test ettiği aralığın dışındadır. Silinmez ama ``ölçülmüş potens değil, eğri ekstrapolasyonu'' olarak etiketlenir.
\item \textbf{Stereokimya korunur.} Koku tanımlayıcıları yalnızca \textbf{tam} InChIKey eşleşmesiyle atanır; farklı bir stereoizomerin etiketi asla ikame edilmez. Gerekçe veride görünür: (R)- ve (S)-karvon farklı reseptörleri etkinleştirir.
\item \textbf{Doğal haplotipler ayrı tutulur.} OR7D4 WM gibi alleller gene katlanmaz, ayrı allel olarak saklanır.
\item \textbf{Eşlenemeyen 4 isim} (OR11H13, OR2T32, ORM2B, ORN4C) olduğu gibi korunur ve özet raporda listelenir.
\end{itemize}


\vspace{0.4\baselineskip}
\hrule
\vspace{0.9\baselineskip}


\section{Simülasyon modeli}

Modelin tamamı React arayüzünden bağımsız, saf TypeScript olarak \texttt{\small src/model/} altında yazılmıştır ve birim testleri vardır. Yani simülasyon arayüzden ayrı olarak test edilebilir ve yeniden kullanılabilir.


\subsection{Reseptör aktivasyonu — \texttt{\small measured}}

Ölçülmüş EC50 varsa, Hill katsayısı 1 olan basit doygunluk eğrisi:


\begin{Verbatim}[fontsize=\small, xleftmargin=1em]
a = C / (EC50 + C)
\end{Verbatim}

Hill katsayısı 1 alınır, çünkü M2OR Hill katsayısı bildirmez — yani bu bir varsayım değil, veri yokluğunun dürüst sonucudur.

Yalnızca taramadan bilinen agonistler, yanıt veren en düşük konsantrasyondan itibaren sabit bir niteliksel düzeye (\textbf{0,5}) geçer. Test aralığının üzerine uydurulmuş bir EC50 sansürlü kabul edilir ve \textbf{0,25} düzeyi verilir.

Arayüzdeki ``barkod'' görünümünde \textbf{işaret boyutu kanıt düzeyini kodlar}: doz–yanıt > doğrulanmış tarama > tek tarama vuruşu. Böylece bir örüntüye bakan kişi, 54 agonistin kaçının gerçekte doz–yanıt eğrisiyle ölçüldüğünü tek bakışta görür.


\subsection{Duysal nöron — \texttt{\small schematic}}

Beş bağlaşık denklemle tek bir olfaktör duysal nöronun kinetik modeli:

reseptör sürüşü \(\rightarrow\) G\(\alpha\)olf ve AC III (\(\tau\) = 25 ms) \(\rightarrow\) cAMP (sentez \texttt{\small kSyn} = 0,02, PDE1C yıkımı \texttt{\small kDeg} = 0,01) \(\rightarrow\) CNG kanal açılması (Hill katsayısı \textbf{2}, Ca\(^{2+}\)–kalmodulin ile duyarlılık azalması, \(\tau\) = 250 ms) \(\rightarrow\) Ca\(^{2+}\) (NCKX4 ile atılım, \(\tau\) = 100 ms) \(\rightarrow\) ANO2 klorür akımı (reseptör akımının \textbf{\%90}'ına kadarını taşır, \texttt{\small gCl} = 9) \(\rightarrow\) depolarizasyon (\(\tau\) = 10 ms) \(\rightarrow\) \textbf{200 Hz}'e kadar ateşleme.

Modelin \textbf{şekilleri} — gecikme, amplifikasyon, adaptasyon, sonlanma — Reisert ve ark.'nın kayıtlarını izler (\href{https://doi.org/10.1085/jgp.201110645}{doi:\allowbreak 10.\allowbreak 1085/\allowbreak jgp.\allowbreak 201110645}). \textbf{Parametre değerleri açıklayıcıdır ve hiçbir kayda fit edilmemiştir.} Zamanlama ve transdüksiyon ayrıntıları büyük ölçüde kemirgen ve amfibi nöronlarından gelir; bu, dokümante edilmiş bir sınırlamadır.


\subsection{Bulbus olfactorius — \texttt{\small schematic}}

Her reseptör tipi için bir glomerül, o nöron tipinin ortalama ateşleme hızını alır; lateral inhibisyon ağırlıklı bir komşuluk çıkarımı yaparak örüntüyü keskinleştirir.

\textbf{Yerleşim anatomik değildir} ve site bunu açıkça söyler: hangi insan reseptörünün hangi glomerüle karşılık geldiği bilinmiyor, ayrıca insanda bulbus başına 5.500'den fazla glomerül vardır — reseptör tipi başına kabaca 16 (\href{https://doi.org/10.1371/journal.pone.0002640}{doi:\allowbreak 10.\allowbreak 1371/\allowbreak journal.\allowbreak pone.\allowbreak 0002640}).


\subsection{Korteks ve algı — \texttt{\small schematic}}

Bölge parlaklığı toplam bulbus çıktısını izler (yarı-doygunluk 0,5); algılama, toplam nöron sürüşünü 0,05 eşiğiyle karşılaştırır. Bunlar açıklayıcı çıktılardır, ölçülmüş aktivite değildir.

\textbf{Önemli:} algı olarak gösterilen koku sözcükleri \textbf{her zaman insan algı veri setlerinden} gelir; animasyondan hesaplanmaz. Yani simülasyon ``vanilya kokuyor'' sonucunu üretmez — o kelime Leffingwell/Dravnieks/Keller–Vosshall verisinden okunur.


\subsection{``Kim kokluyor'' — birey değişkenleri}

Tüm çarpanlar tek bir dosyada (\texttt{\small src/model/modifiers.ts}), her biri kaynağı ve gerekçe yorumuyla birlikte toplanmıştır. İki ilke:

\begin{itemize}
\item \textbf{Yalnızca reseptör genotipi, reseptör aktivasyonunu değiştirir} — ve bunu yaparken \textbf{tek bir makale içinde ölçülmüş} allel verisini kullanır. Böylece genotip etkisi laboratuvarlar arası farklarla karışmaz. Heterozigot, her allelden yarım olarak modellenir; gerekçe monoallelik ifadedir (kemirgen verisi, \href{https://doi.org/10.1016/S0092-8674(94}{doi:\allowbreak 10.\allowbreak 1016/\allowbreak S0092-8674(94})90562-2).
\item \textbf{Yaş}, reseptör aşamasından \textbf{sonra} tek bir küçük ve açıkça ``illustrative'' olarak etiketlenmiş kaydırma uygular (0,82 çarpanı; primer olfaktör korteks, insula, hipokampus). Çalışmadan gelen tek şey \textbf{yön}dür (yaşlı yetişkinlerde daha düşük), çarpanın büyüklüğü değil (fMRI, n = 37; \href{https://doi.org/10.18632/oncotarget.16977}{doi:\allowbreak 10.\allowbreak 18632/\allowbreak oncotarget.\allowbreak 16977}).
\end{itemize}

Genotip vakaları:


\begin{small}
\begin{longtable}{@{}p{0.149\linewidth}@{\hspace{5pt}}@{}p{0.106\linewidth}@{\hspace{5pt}}@{}p{0.167\linewidth}@{\hspace{5pt}}@{}p{0.296\linewidth}@{\hspace{5pt}}@{}p{0.211\linewidth}@{}}
\toprule
\textbf{Vaka} & \textbf{Reseptör} & \textbf{Odorant} & \textbf{Alleller} & \textbf{Kaynak} \\
\midrule
\endfirsthead
\toprule
\textbf{Vaka} & \textbf{Reseptör} & \textbf{Odorant} & \textbf{Alleller} & \textbf{Kaynak} \\
\midrule
\endhead
\bottomrule
\endfoot
Androstenon & OR7D4 & androstenon & RT (referans), WM (R88W+T133M) & \href{https://doi.org/10.1038/nature06162}{doi:\allowbreak 10.\allowbreak 1038/\allowbreak nature06162} \\
cis-3-heksenol & OR2J3 & cis-3-heksen-1-ol & referans, T113A+R226Q & \href{https://doi.org/10.1093/chemse/bjs049}{doi:\allowbreak 10.\allowbreak 1093/\allowbreak chemse/\allowbreak bjs049} \\
Guaiakol & OR10G4 & guaiakol & referans, APT (2 aa), VLT (5 aa), ALI (8 aa) & \href{https://doi.org/10.1038/nn.3598}{doi:\allowbreak 10.\allowbreak 1038/\allowbreak nn.\allowbreak 3598} \\
\end{longtable}
\end{small}


\vspace{0.4\baselineskip}
\hrule
\vspace{0.9\baselineskip}


\section{Ölçülmemiş bir molekülü tahmin etmek — \texttt{\small predicted}}

``Molekül çiz'' ekranında kullanıcı herhangi bir yapıyı çizebilir veya SMILES olarak yapıştırabilir; sistem reseptör örüntüsünü ve koku sözcüklerini tahmin eder. Tahmin aynı \texttt{\small Stimulus} arayüzünü beslediği için \textbf{tüm simülasyon çizilmiş bir molekülle de değişmeden çalışır}, baştan sona \texttt{\small predicted} etiketiyle.


\subsection{Moleküler temsil}

Morgan parmak izi, yarıçap 2, 2048 bit, \textbf{kiralite dahil}. Kiralite dahil edilmesi bilinçli bir karardır: ölçülmüş veri enantiyomerlerin farklı reseptörleri etkinleştirdiğini gösteriyor, dolayısıyla (R)- ve (S)-karvonu aynı molekül sayan bir tahminci sitenin kendi verisiyle çelişirdi.

Python (RDKit) ve tarayıcı (RDKit.js) \textbf{bit düzeyinde aynı} parmak izini üretir; bu nedenle referans seti önceden hesaplanır ve çalışma zamanında yalnızca sorgu molekülü parmak izlenir.


\subsection{Yöntem A — en yakın komşular (temel çizgi)}

M2OR'dan \textbf{654 ölçülmüş molekül}e karşı Tanimoto benzerliği. Bir reseptörün skoru, benzer moleküllerin yanıt verenlerinin benzerlik-ağırlıklı payıdır — \textbf{yalnızca o reseptörde gerçekten test edilmiş olanlar sayılarak.} Böylece hiç test edilmemiş bir reseptör ``negatif tahmin'' edilmez, \textbf{bilinmeyen} kalır. Bir reseptör ancak toplam benzerliği \(\geq\) 1,0 olan komşular onda test edilmişse listelenir ve skor > 0,5 ise yanıt veren sayılır. Koku sözcükleri aynı yolla, uzman etiketli \textbf{3.522 molekül}den gelir. Her tahmin, hangi ölçülmüş moleküllerden türediğini listeler.

\textbf{Birini-dışarıda-bırak değerlendirmesi} (142 molekül skorlandı, yalnızca gerçekten test edildikleri reseptörler puanlanarak):


\begin{small}
\begin{longtable}{@{}p{0.578\linewidth}@{\hspace{5pt}}@{}p{0.352\linewidth}@{}}
\toprule
\textbf{Metrik} & \textbf{Değer} \\
\midrule
\endfirsthead
\toprule
\textbf{Metrik} & \textbf{Değer} \\
\midrule
\endhead
\bottomrule
\endfoot
En üst sıradaki reseptör gerçek agonist & \textbf{\%61} (n = 138) \\
Kesinlik (precision) @ skor \(\geq\) 0,5 & 0,48 \\
Duyarlılık (recall) @ skor \(\geq\) 0,5 & 0,13 \\
\end{longtable}
\end{small}

\textbf{Bu sayıların dürüst okuması:} sıralamayı okuyun, mutlak kümeyi değil. Duyarlılığın düşük olmasının yapısal bir nedeni var — her referans molekülü yalnızca birkaç reseptörde test edilmişken, ``gerçeği'' tanımlayan taramalar yüzlerce reseptörü taradı. Tahmin listesinin tepesi, \textbf{test edilmeye değer bir hipotez}dir; bir sonuç değildir.


\subsection{Yöntem B — öğrenilmiş model}

Her iki yanda da yeterli ölçülmüş molekülü olan \textbf{31 reseptör} için (her sınıfta \(\geq\) 8 molekül), parmak izi bitleri üzerinde reseptör başına \textbf{L2 düzenlileştirilmiş lojistik regresyon}. Komşuluk temel çizgisine karşı \textbf{birebir aynı katlarda} çapraz doğrulanır:


\begin{small}
\begin{longtable}{@{}p{0.392\linewidth}@{\hspace{5pt}}@{}p{0.264\linewidth}@{\hspace{5pt}}@{}p{0.274\linewidth}@{}}
\toprule
\textbf{31 reseptör üzerinde} & \textbf{Öğrenilmiş model} & \textbf{En yakın komşular} \\
\midrule
\endfirsthead
\toprule
\textbf{31 reseptör üzerinde} & \textbf{Öğrenilmiş model} & \textbf{En yakın komşular} \\
\midrule
\endhead
\bottomrule
\endfoot
Ortalama ROC-AUC & \textbf{0,691} & 0,677 \\
Ortalama average precision & \textbf{0,422} & 0,380 \\
Daha iyi olduğu reseptör sayısı & 22 & 9 \\
\end{longtable}
\end{small}

Fark marjinaldir ve \textbf{abartılmamalıdır.}

Asıl değişken \textbf{reseptörün kendisidir}: en iyi model OR5AN1 için AUC 0,98'e ulaşırken birkaç reseptör şans düzeyinde kalır. Tek bir manşet sayısı bunu gizlerdi; bu yüzden \textbf{her tahmin kendi reseptörünün çapraz doğrulanmış AUC'sini taşır}, 0,70 altındaki değerler güvenilmez olarak işaretlenir ve simülasyon öğrenilmiş skoru \textbf{yalnızca bu çizginin üstünde} kullanır. Model 43 KB'lık bir matrise aktarılır ve tarayıcıda basit bir iç çarpım olarak çalışır.


\vspace{0.4\baselineskip}
\hrule
\vspace{0.9\baselineskip}


\section{Benchmark: düz çapraz doğrulama ile iç içe çapraz doğrulama}

Yukarıdaki tablo \emph{bu model ne kadar iyi} sorusunu yanıtlar. \emph{Bu modeli seçen prosedür ne kadar iyi} sorusunu yanıtlamaz — ve bir şeyi ayarlamaya (tuning) başladığınız anda önemli olan sayı ikincisidir, çünkü \textbf{bir kazananı seçen katlar, aynı zamanda onu adil notlandıramaz.}

\texttt{\small scripts/pipeline/benchmark\_models.py} dokuz model ailesini birebir aynı katlarda iki şemadan geçirir (5-katlı tabakalı dış \(\times\) 3 karıştırma, havuzlanmış kat-dışı skorlama, 3-katlı iç arama, bit filtreleme her katta yeniden fit edilir; \textasciitilde{}24 dakika):


\begin{small}
\begin{longtable}{@{}p{0.264\linewidth}@{\hspace{5pt}}@{}p{0.083\linewidth}@{\hspace{5pt}}@{}p{0.083\linewidth}@{\hspace{5pt}}@{}p{0.106\linewidth}@{\hspace{5pt}}@{}p{0.113\linewidth}@{\hspace{5pt}}@{}p{0.182\linewidth}@{\hspace{5pt}}@{}p{0.099\linewidth}@{}}
\toprule
\textbf{Model ailesi} & \textbf{Izgara} & \textbf{Düz CV} & \textbf{İç içe CV} & \textbf{İyimserlik} & \textbf{\(\Delta\) (sevk edilene karşı)} & \textbf{p (Holm)} \\
\midrule
\endfirsthead
\toprule
\textbf{Model ailesi} & \textbf{Izgara} & \textbf{Düz CV} & \textbf{İç içe CV} & \textbf{İyimserlik} & \textbf{\(\Delta\) (sevk edilene karşı)} & \textbf{p (Holm)} \\
\midrule
\endhead
\bottomrule
\endfoot
lojistik C=0,1 (\textbf{sevk edilen}, ayarsız) & 1 & 0,686 & \textbf{0,686} & — & — & — \\
random forest & 4 & 0,707 & 0,686 & +0,022 & 0,000 & 0,78 \\
extra trees & 4 & 0,710 & 0,683 & +0,027 & --0,003 & 0,78 \\
lojistik regresyon (L2) & 6 & 0,707 & 0,677 & +0,030 & --0,009 & 0,26 \\
Bernoulli naive Bayes & 3 & 0,701 & 0,676 & +0,025 & --0,010 & 0,52 \\
Tanimoto komşuları & 3 & 0,703 & 0,674 & +0,029 & --0,012 & 0,52 \\
doğrusal SVM & 4 & 0,712 & 0,671 & +0,040 & --0,015 & 0,08 \\
RBF SVM & 9 & 0,726 & 0,668 & +0,058 & --0,018 & 0,08 \\
lojistik regresyon (L1) & 5 & 0,709 & 0,657 & +0,051 & --0,029 & 0,20 \\
\end{longtable}
\end{small}

Üç sonuç ve hepsi aynı yönü gösteriyor:

\begin{enumerate}
\item \textbf{Hiçbir şey, ayarsız sevk edilen yapılandırmayı geçmiyor.} Dokuz ailenin en yüksek iç içe CV skoru odur (random forest ile eşit). Sekiz karşılaştırma üzerinde Holm düzeltmesinden sonra hiçbir aile ondan anlamlı biçimde farklı değil; düzeltilmemiş \emph{p} < 0,05'e ulaşan üç aile ise \textbf{daha kötü} olanlardır.
\item \textbf{Ayarlama, var olmayan +0,035 AUC değerinde.} Bu, ortalama düz-eksi-iç içe farkıdır — hiperparametre aramasının bedeli dürüstçe ödendiğinde buharlaşan skor. Bu fark, dokuz aile arasındaki tüm yayılımdan daha büyüktür; yani düz CV ile raporlanan ayarlı bir model, hiç olmadığı hâlde iyileşme gibi görünürdü. Fark ızgara boyutuyla ölçeklenir: dokuz noktalı RBF SVM ızgarası için +0,058, hiçbir şey aramayan model için +0,000.
\item \textbf{Kat gürültüsü, farklarla aynı büyüklükte.} Aynı değerlendirme farklı bir karıştırmayla yeniden koşulduğunda tek bir reseptörün AUC'si ±0,021 oynuyor.
\end{enumerate}

\textbf{Sonuç:} bu veri setinde sınır algoritma değil, \textbf{veridir.} Bir sonraki kazanç daha geniş bir aramadan değil, ya daha fazla ölçülmüş çiftten ya da reseptörler arasında tek bir temsili paylaşmaktan gelmek zorunda. Karşılaştırma gizlenmek yerine \textbf{yayımlanmıştır} — sitede \texttt{\small \#/benchmark} sayfası tabloyu, düz-karşı-iç içe şeklini ve her reseptörün kendi sayılarını gösterir; \texttt{\small scripts/report/benchmark\_report.py} aynı JSON'dan tek sayfalık bir LaTeX özeti üretir, böylece ikisi birbirinden sapamaz.


\vspace{0.4\baselineskip}
\hrule
\vspace{0.9\baselineskip}


\section{Karışımlar}

Karışımlar ekranında bileşenler her reseptörde \textbf{aynı bağlanma bölgesi için yarışır}. Eşit etkinlikte agonistler için, C\(_{i}\) konsantrasyonları ve EC50\(_{i}\) potensleriyle:


\begin{Verbatim}[fontsize=\small, xleftmargin=1em]
a = Sum(Ci/EC50i) / (1 + Sum(Ci/EC50i))        Ci = fi*C yazılırsa:
a = C / (EC50_karışım + C)                  EC50_karışım = 1 / Sum(fi/EC50i)
\end{Verbatim}

Yani karışım, reseptör başına \textbf{tek bir etkin EC50}'ye çöker — sıradan bir Hill eğrisi — ve tek bir \texttt{\small Stimulus} hâline gelir; sonraki tüm aşamalar hiç değişmeden çalışır.

Ekranda görülmeye değer iki sonuç:

\begin{itemize}
\item Sabit \textbf{toplam} konsantrasyonda bir karışım, en iyi tek bileşeninden asla daha güçlü değildir; çünkü zayıf bir ortak eklemek güçlü olanı seyreltir.
\item Bir bileşenin bir reseptörün sürüşündeki payı, karışımdaki payını değil \textbf{o reseptördeki potensini} izler. İki karvon enantiyomeri ve guaiakolün 1:1:1 karışımında (R) formu OR1A1'deki sürüşün \textasciitilde{}\%90'ını sağlarken, (S) formu OR2W1'de \textasciitilde{}\%80'ini sağlar.
\end{itemize}

\textbf{Modelin açık sınırı:} bu formülasyon bileşen eklendikçe bir reseptörün yanıtının yalnızca \textbf{artmasına} izin verir. Gerçek karışımlarda görülen \textbf{supresyon ve maskeleme} (reseptör düzeyinde antagonizma ve daha ileri aşamalarda) modellenmemiştir ve bir karışım hiçbir bileşeninin kokmadığı bir şey gibi kokabilir. Ayrıca karışım için gösterilen harmanlanmış koku sözcükleri, bileşenlerin kendi bildirilen sözcüklerinin ağırlıklı bir \textbf{gösterimidir}, karışımın nasıl algılandığına dair bir tahmin değildir. Site bunu belirgin biçimde söyler.


\vspace{0.4\baselineskip}
\hrule
\vspace{0.9\baselineskip}


\section{Sınırlar — projenin iddia \textbf{etmediği} şeyler}

Bu bölüm, nörobilimci bir dinleyiciyle konuşurken en önemli bölümdür ve sitede ayrı bir başlık olarak durur:

\begin{itemize}
\item Reseptör–odorant verisi \textbf{hücre temelli tahlillerden} gelir (HEK293T, Hana3A ve diğerleri; lusiferaz, cAMP veya Ca\(^{2+}\) okumalarıyla). Bunlar hangi reseptörlerin yanıt \textbf{verebileceğini} gösterir, canlı bir burunda ne olduğunu göstermez.
\item İnsan reseptörlerinin çoğunun bilinen ligandı yoktur. \textbf{``Test edilmemiş'', ``yanıt vermiyor'' değildir}; site bunları ayrı gösterir.
\item Tahlil konsantrasyonları, havadaki veya nazal mukustaki konsantrasyonlar değildir.
\item Transdüksiyon ayrıntıları ve zamanlamaları çoğunlukla kemirgen ve amfibi nöronlarındandır.
\item Glomerüler harita, bölge aktivite düzeyleri ve reseptör aşamasından sonraki bireysel kaydırmalar şematik veya açıklayıcıdır.
\item Mainland ve ark. 2015 EC50 değerleri tam-log değerlerdir (±0,5 log) ve işaretlenir.
\item OR51E2 için dokuz EC50 değeri, tahlillerin test ettiği aralığın altına (1 nM altı) uydurulmuştur; korunur ama işaretlenir.
\item Tahminler bir benzerlik temel çizgisi ve küçük bir doğrusal modeldir: sıralama için yararlı, duyarlılıkta zayıf.
\item Karışım modeli supresyonu modellemez (yukarıya bakınız).
\end{itemize}


\vspace{0.4\baselineskip}
\hrule
\vspace{0.9\baselineskip}


\section{Çıktılar, yeniden üretilebilirlik ve atıf}


\subsection{Veri dışa aktarımı}

Reseptör örüntüsü hesaplayan her görünüm, sonucu \textbf{CSV} veya \textbf{JSON} olarak sunar (\texttt{\small src/model/export.ts}). İkisi de 432 reseptörün tamamını, konsantrasyonu, InChIKey veya karışım bileşimini, genotipi ve her değerin kaynağını taşır.

Kimsenin test etmediği reseptörler, \texttt{\small 0} değil \textbf{boş hücre} (CSV) veya \texttt{\small null} (JSON) olarak aktarılır. Böylece ``veri yok'', aşağı akış analizinde sessizce ``yanıt yok''a dönüşemez — arayüzün yaptığı ayrımın aynısı.


\subsection{Paylaşılabilir bağlantılar}

Odorant, konsantrasyon, genotip veya karışım URL hash'inde kodlanır (\texttt{\small src/state/shareLink.ts}), böylece tam bir görünüm bir meslektaşa gönderilebilir veya bir makalede atıf olarak verilebilir:


\begin{Verbatim}[fontsize=\small, xleftmargin=1em]
#/simulate?o=androstenone&c=-5.00&g=OR7D4-androstenone:WM/WM
#/mixtures?c=-4.00&m=r-carvone:1&m=guaiacol:1
\end{Verbatim}

Bilinmeyen odorantlar, aralık dışı konsantrasyonlar ve bozuk değerler hata vermek yerine yok sayılır — çünkü bu bağlantılar elle yeniden yazılır ve kırpılır.


\subsection{Atıf}

Site, \textbf{başkalarının ölçümlerinin bir sunumudur.} Bir şekil veya dışa aktarılmış örüntü kullanılırsa \textbf{veri kaynaklarına atıf verilmelidir} — iş onların üzerinde duruyor:

\begin{itemize}
\item \textbf{Reseptör–odorant verisi:} Lalis, M., Hladiš, M., Abi Khalil, S., Briand, L., Fiorucci, S. \& Topin, J. M2OR: a database of olfactory receptor–odorant pairs. \emph{Nucleic Acids Research} (2024). \href{https://doi.org/10.1093/nar/gkad886}{doi:\allowbreak 10.\allowbreak 1093/\allowbreak nar/\allowbreak gkad886}
\item \textbf{Koku tanımlayıcıları:} Leffingwell PMP 2001 (Pyrfume üzerinden, CC BY-NC), Dravnieks \emph{Atlas of Odor Character Profiles} (1985), Keller \& Vosshall (\href{https://doi.org/10.1186/s12868-016-0287-2}{doi:\allowbreak 10.\allowbreak 1186/\allowbreak s12868-016-0287-2})
\item \textbf{Reseptör kimliği ve yapılar:} HGNC, UniProt, GPCRdb, AlphaFold DB ve her PDB girdisinin birincil atfı
\item \textbf{Araç için:} ``Molekülden kokuya'', interaktif olfaktör reseptör simülatörü
\end{itemize}

Hangi genotip, yapı veya EC50 alıntılanıyorsa, o değerin yanında \textbf{spesifik makale} bağlantılıdır — ölçümün kendisi için siteye değil o makaleye atıf verilmelidir.


\vspace{0.4\baselineskip}
\hrule
\vspace{0.9\baselineskip}


\section{Teknik mimari}


\begin{Verbatim}[fontsize=\small, xleftmargin=1em]
scripts/pipeline/   Python veri hattı (pandas, RDKit, gemmi) + model benchmark'ı
scripts/report/     Benchmark JSON'undan üretilen LaTeX raporu
reports/            Üretilmiş tek sayfalık rapor (.tex ve .pdf)
data/raw/           İndirmeler (git tarafından yok sayılır, önbelleklenir)
data/processed/     Sitenin yüklediği JSON (depoya işlenmiş)
src/model/          Saf TypeScript simülasyon, React yok — birim testli
src/scenes/         Aşama başına bir SVG sahne
src/ui/             Yeniden kullanılabilir bileşenler (barkod,
                    grafikler, moleküller, 3B, kontroller)
src/screens/        Başlangıç, Simülasyon, Reseptör gezgini,
                    Odorant gezgini, Molekül çiz, Karışımlar,
                    Nasıl çalışır
attic/              Derlemeden çıkarılmış eski dosyalar
\end{Verbatim}

\texttt{\small src/model/} arayüzden \textbf{hiç import etmez.} Modüller:


\begin{small}
\begin{longtable}{@{}p{0.284\linewidth}@{\hspace{5pt}}@{}p{0.646\linewidth}@{}}
\toprule
\textbf{Dosya} & \textbf{İçerik} \\
\midrule
\endfirsthead
\toprule
\textbf{Dosya} & \textbf{İçerik} \\
\midrule
\endhead
\bottomrule
\endfoot
\texttt{\small activation.ts} & Hill fonksiyonu, yalnız-tarama agonistleri için niteliksel düzeyler \\
\texttt{\small transduction.ts} & Tek duysal nöronun kinetik modeli (cAMP, Ca\(^{2+}\), Cl\(^{-}\), Vm, ateşleme) \\
\texttt{\small bulb.ts} & Glomerüler yerleşim ve lateral inhibisyon \\
\texttt{\small simulation.ts} & Tam bir koşu \\
\texttt{\small modifiers.ts} & Her birey-seçici çarpanı, kaynağı ve yorumuyla \\
\texttt{\small stimulus.ts} & Faz 2 ve 3'ün bağlandığı arayüz (seam) \\
\end{longtable}
\end{small}


\subsection{Çalıştırma}


\begin{Verbatim}[fontsize=\small, xleftmargin=1em]
npm install          # bağımlılıklar
npm run dev          # http://localhost:5173
npm run build        # tip kontrolü + dist/ içine üretim derlemesi
npm test             # model birim testleri (Vitest)
\end{Verbatim}

\texttt{\small npm run dev}, \texttt{\small data/processed/*.json} ile Mol*/RDKit paketlerini önce \texttt{\small public/} içine kopyalar; yani \textbf{depodaki veri siteyi çalıştırmak için yeterlidir.}


\subsection{Veriyi yeniden üretme}


\begin{Verbatim}[fontsize=\small, xleftmargin=1em]
python3 -m venv .venv
.venv/bin/pip install -r scripts/requirements.txt
.venv/bin/python -m scripts.pipeline.build            # önbelleklenmiş indirmelerle
.venv/bin/python -m scripts.pipeline.build --refresh  # kaynakları yeniden indir
npm run benchmark    # ~24 dk -> data/processed/model_benchmark.json
npm run report       # -> reports/benchmark_report.pdf
\end{Verbatim}

\textbf{Tek manuel adım.} M2OR web sitesi (m2or.chemsensim.fr) betiklerin geçemediği bir tarayıcı kontrolünün arkasındadır, bu nedenle hat yazarların GitHub anlık görüntüsünü kullanır (\texttt{\small M2OR\_20230428.csv}, Apache-2.0). Güncel veritabanını kullanmak için Experiments tablosunu tarayıcıdan CSV olarak dışa aktarıp \texttt{\small data/raw/m2or/} içine bırakmak yeterlidir; hat onu otomatik olarak tercih eder ve farkları özet raporda bildirir.


\subsection{Yayımlama}

Derleme, göreli tabanlı, kendi kendine yeten statik bir sitedir — sunucu, veritabanı veya API anahtarı yok. \texttt{\small dist/} yaklaşık 21 MB'dır (3B görüntüleyici, kimya araç takımı ve çizim aracı ağırlığın çoğunu oluşturur ve yalnızca gerektiğinde yüklenir). \texttt{\small .github/workflows/deploy.yml}, \texttt{\small main}'e her itişte derleyip GitHub Pages'e yayımlar.


\vspace{0.4\baselineskip}
\hrule
\vspace{0.9\baselineskip}


\section{Fazlar ve yol haritası}

\begin{itemize}
\item \textbf{Faz 1 (hazır):} bir test odorantı seç, tam simülasyonu izle, her reseptörü incele.
\item \textbf{Faz 2 (hazır):} herhangi bir molekülü çiz veya yapıştır; tahmini reseptör örüntüsü ve tahmini koku sözcüklerini, her tahminin arkasındaki ölçülmüş molekülleri kanıt olarak göstererek al.
\item \textbf{Faz 3 (hazır):} iki veya daha fazla odorantın (çizilmiş bir molekül dahil) karışımı, aynı \texttt{\small Stimulus}'a toplanarak.
\end{itemize}

\textbf{Planlanan değişiklik:} aynı arayüzün arkasına daha yüksek kapasiteli bir öğrenilmiş model — reseptörler için protein dil modeli ile bir graf sinir ağı (Hladiš ve ark. 2023) ve koku tanımlayıcıları için OpenPOM. Ancak yukarıdaki kanıta göre \textbf{daha fazla veri, daha iyi bir mimariden daha çok işe yarayacaktır.}


\subsection{İlgili yeni literatür}

Pashakhanloo \& Murthy (Harvard, arXiv:2610.17164, Ekim 2026), \textasciitilde{}1.240 doğal gıda karışımının moleküler bileşiminden algısal benzerliğini öğrenmek için büyük dil modellerinden türetilen koku tanımlayıcılarını vekil hedef olarak kullanıyor. Bu proje açısından üç noktası ilgi çekici:

\begin{enumerate}
\item Doğrulamayla seçilen moleküler gömme boyutu \textbf{k \(\approx\) 141} (\%95 GA 99–200) çıkıyor; yazarlar bunun insan OR tipi sayısıyla aynı mertebede olduğunu ve reseptör repertuvarının \textbf{biyolojik bir darboğaz} oluşturduğunu vurguluyor. Hiç reseptör verisi kullanılmadan, yalnızca algısal benzerlik uydurmasından bu projenin merkezî tezine bağımsız destek.
\item Dışsal moleküler öznitelikler (GS–LF, OpenPOM, MolFormer, Mordred) sıfır-vuruşta yalnızca r \(\approx\) 0,24–0,35 veriyor ve \textbf{yalnızca az veri rejiminde} fayda sağlıyor; tam eğitim setinde ek katkıları kayboluyor. Bu, bu projenin benchmark sonucunun (``sınır algoritma değil, veri'') bağımsız bir doğrulamasıdır.
\item Yazarların açık sınırlaması, VCF verisinde güvenilir konsantrasyon olmadığı için karışımları \textbf{ikili var/yok} vektörleriyle temsil etmek zorunda kalmaları. Bu projenin karışım modeli ise doğrudan konsantrasyon sürümlüdür.
\end{enumerate}


\vspace{0.4\baselineskip}
\hrule
\vspace{0.9\baselineskip}


\section{Toplantıda muhtemel sorular ve kısa yanıtlar}

\textbf{``Bu bir tahmin motoru mu, bir görselleştirme mi?''} Reseptör adımında bir veri sunumu; ``Molekül çiz'' ekranında ise açıkça etiketlenmiş bir tahmin. İkisi hiçbir yerde karışmaz, çünkü her değer provenance etiketi taşır.

\textbf{``Neden Hill katsayısı 1?''} Çünkü M2OR Hill katsayısı bildirmiyor. Veri olmadığında parametre uydurmak yerine en yalın varsayım seçildi ve bu açıkça yazıldı.

\textbf{``235 reseptörün agonisti var — gerisi yanıt vermiyor mu?''} Hayır. 422 reseptör en az bir molekülle test edilmiş; 235'inde en az bir agonist bulundu. Kalanı için çoğunlukla \textbf{yeterli test yok.} Site ve dışa aktarım bu iki durumu ayırır.

\textbf{``Bulbus haritası gerçek mi?''} Değil ve site bunu söylüyor. İnsan reseptör\(\rightarrow\)glomerül eşlemesi bilinmiyor; ayrıca insanda reseptör tipi başına \textasciitilde{}16 glomerül var. Yerleşim açıklayıcıdır.

\textbf{``Tahmin performansı düşük görünüyor (recall 0,13).''} Doğru ve bu yapısal. Referans molekülleri birkaç reseptörde test edilmişken ``gerçeği'' tanımlayan taramalar yüzlerce reseptör taradı. Bu yüzden mutlak kümeye değil \textbf{sıralamaya} bakılması söyleniyor ve her tahmin kendi reseptörünün AUC'siyle sunuluyor.

\textbf{``Daha iyi bir model denediniz mi?''} Dokuz model ailesi, iç içe çapraz doğrulamayla. Hiçbiri sevk edilen ayarsız yapılandırmayı güvenilir biçimde geçmedi; ayarlama düz CV'de \textasciitilde{}0,03 AUC kazanç gibi görünüp iç içe şemada buharlaşıyor. Karşılaştırma sitede yayımlanmış durumda.

\textbf{``Karışımlarda supresyon neden yok?''} Çünkü mevcut veri yarışmalı bağlanmanın ötesini kalibre etmeye yetmiyor. Model yalnızca tek bağlanma bölgesi için yarışmayı içeriyor ve supresyonun modellenmediği açıkça belirtiliyor.


\end{document}
